% A line starting with a % character is a comment and wont be included in the report.
% Always use two backslashes like \\ to insert a new line. Take a look at the example below.
For multiple target tracking, I implemented a system using multiple Kalman Filters (one per target) rather than Particle Filters. Here's what I did and why: \\

\textbf{Implementation Approach:} \\
1. \textbf{Multiple Kalman Filters}: Created separate KF instances for each detected person (up to 3 targets) \\
2. \textbf{HOG Detection}: Used Histogram of Oriented Gradients (HOG) detector to identify people in each frame \\
3. \textbf{Target Association}: Associated detections with existing trackers based on proximity to previous positions \\
4. \textbf{Independent Tracking}: Each target was tracked independently with its own state and covariance \\

\textbf{KF vs PF for Multiple Targets:} \\

\textbf{Kalman Filter Advantages:} \\
- \textbf{Computational Efficiency}: Much faster than particle filters, especially with multiple targets \\
- \textbf{Deterministic}: Predictable behavior, easier to debug \\
- \textbf{Linear Dynamics}: Well-suited for smooth motion prediction \\
- \textbf{Memory Efficient}: Only stores mean and covariance per target \\

\textbf{Particle Filter Disadvantages for Multiple Targets:} \\
- \textbf{Computational Cost}: Would require 100-300 particles per target, leading to thousands of particles total \\
- \textbf{Memory Intensive}: Each particle stores full state information \\
- \textbf{Resampling Complexity}: More complex resampling when targets interact \\

\textbf{Why KF Worked Best:} \\
1. \textbf{Scalability}: Can easily handle 3+ targets without performance degradation \\
2. \textbf{Real-time Performance}: Fast enough for video processing \\
3. \textbf{Robust Association}: Simple distance-based association works well with KF predictions \\
4. \textbf{Smooth Tracking}: Linear motion model provides smooth trajectories \\

\textbf{Modifications for Multiple Targets:} \\
- \textbf{Target Initialization}: Only initialize trackers on first frame with strongest detections \\
- \textbf{Association Logic}: Match detections to trackers based on minimum distance \\
- \textbf{Independent States}: Each tracker maintains separate state and covariance \\
- \textbf{Visualization}: Draw separate bounding boxes and labels for each target \\

The KF approach proved superior for multiple target tracking due to its computational efficiency and deterministic behavior, making it ideal for real-time applications with multiple moving objects.
